%------------------------------------------------------------------------------
% Chapter 1b: The 10C CORE Architecture
%------------------------------------------------------------------------------
% Version: 2.0 (January 2026)
%
% Purpose: Introduces the complete 10C CORE Framework---the eight fundamental
%          questions that structure behavioral prediction in EBF. This chapter
%          provides the architectural overview; each CORE has a dedicated appendix.
%
% Primary Appendix: None (Foundation Chapter introducing all 8 CORE appendices)
% Secondary Appendices: AAA (WHO), B (HOW), C (WHAT), V (WHEN), BBB (WHERE),
%                       AU (AWARE), AV (READY), AW (STAGE), G (Glossary)
%
% Prerequisites: Chapter 1 (Introduction)
% Leads to: Chapter 2 (Rationality), Chapter 5 (Complementarity)
%
% Chapter Type: B (Foundation)
%------------------------------------------------------------------------------

\section{The 10C CORE Architecture}
\label{sec:10c-core}

%------------------------------------------------------------------------------
% QUICK REFERENCE BOX
%------------------------------------------------------------------------------
\begin{tcolorbox}[colback=gray!5!white, colframe=gray!75!black,
    title=\textbf{Begriffe in diesem Kapitel}]

\small
\begin{tabular}{@{}p{3cm}p{9.5cm}@{}}
\textbf{10C CORE} & Die acht fundamentalen Fragen: WHO, WHAT, HOW, WHEN, WHERE, AWARE, READY, STAGE \\[3pt]
\textbf{$U_{pot}$} & Potentielle Utility---was theoretisch möglich wäre \\[3pt]
\textbf{$U_{eff}$} & Effektive Utility---was psychologisch salient ist ($U_{eff} = A \cdot U_{pot}$) \\[3pt]
\textbf{WAX} & Willingness-to-Act Index---muss Threshold $\theta$ überschreiten \\[3pt]
\textbf{$S(t)$} & Stage---Position auf der Behavioral Change Journey (0--5) \\[3pt]
\textbf{4-Stage Pipeline} & $U_{pot} \rightarrow U_{eff} \rightarrow WAX \geq \theta \rightarrow S(t)$ \\
\end{tabular}

\vspace{3pt}
\textit{Vollständiges Glossar: Appendix~G}
\end{tcolorbox}

%------------------------------------------------------------------------------
% APPENDIX REFERENCES BOX
%------------------------------------------------------------------------------
\begin{tcolorbox}[
    colback=yellow!10!white,
    colframe=orange!75!black,
    title={\textbf{Appendix References for This Chapter}},
    fonttitle=\bfseries
]
\small
\begin{tabular}{@{}p{3.2cm}p{5.2cm}p{3.8cm}@{}}
\textbf{Appendix} & \textbf{Content} & \textbf{Section} \\
\midrule
AAA (CORE-WHO) & Level hierarchy (9 levels, 41 axioms) & \ref{sec:10c-who} \\
C (CORE-WHAT) & FEPSDE dimensions (144 components) & \ref{sec:10c-what} \\
B (CORE-HOW) & Complementarity framework & \ref{sec:10c-how} \\
V (CORE-WHEN) & Context $\Psi$ (8 dimensions) & \ref{sec:10c-when} \\
BBB (CORE-WHERE) & Calibration methodology & \ref{sec:10c-where} \\
AU (CORE-AWARE) & Awareness filter (104 axioms) & \ref{sec:10c-aware} \\
AV (CORE-READY) & Willingness threshold (99 axioms) & \ref{sec:10c-ready} \\
AW (CORE-STAGE) & Behavioral change journey & \ref{sec:10c-stage} \\
G (REF-GLOSSARY) & Symbol definitions & All sections
\end{tabular}
\end{tcolorbox}

%------------------------------------------------------------------------------
% INTUITION BOX
%------------------------------------------------------------------------------
\begin{tcolorbox}[colback=gray!5!white, colframe=gray!75!black,
    title=\textbf{Why This Matters: Thomas's Fitness Journey}]
\small
Thomas \emph{knows} he should exercise more. His doctor told him (AWARE = high). He even \emph{wants} to---the gym is 5 minutes away (READY: low barriers). Yet months pass without action.

\textbf{Standard economics} says: ``If utility is positive, action follows.'' Thomas has positive utility from exercise. Why no action?

\textbf{The 10C Framework} explains:
\begin{itemize}[nosep]
    \item \textbf{WHO:} Thomas thinks as individual ($L=0$), not family ($L=2$)---his kids' health isn't salient
    \item \textbf{WHAT:} He focuses on Physical ($P$), missing Social ($S$) benefits of gym community
    \item \textbf{AWARE:} Long-term health benefits ($U_{pot}$) filtered by present bias---$U_{eff} \ll U_{pot}$
    \item \textbf{READY:} $WAX = 0.6$ but threshold $\theta = 0.7$---not quite ready
    \item \textbf{STAGE:} $S = 1$ (Contemplation)---thinking about it, not acting
\end{itemize}

\textit{This chapter shows how the 10C Framework predicts exactly WHEN and WHY behavior change occurs---or fails.}
\end{tcolorbox}

%------------------------------------------------------------------------------
% CENTRAL QUESTION BOX
%------------------------------------------------------------------------------
\paragraph{The Central Question.}
This chapter addresses the core architectural question of EBF:

\begin{center}
\fcolorbox{blue!75!black}{blue!5!white}{\parbox{0.85\textwidth}{
\textbf{Central Question:} What are the fundamental questions that must be answered to predict human behavior---and how do they connect into a complete prediction architecture?
}}
\end{center}

%------------------------------------------------------------------------------
% BRIDGE FROM CHAPTER 1
%------------------------------------------------------------------------------
\paragraph{Building on Chapter 1.}
In Chapter~1, we introduced four core concepts: \textbf{Complementarity} ($C$), \textbf{Context} ($\Psi$), the \textbf{Reference Structure} ($C^*$), and \textbf{Coherence} ($K$). These concepts explained \emph{why} Anna in Switzerland and Anna in Peru behave differently in the ultimatum game---not because they have different preferences, but because context ($\Psi$) modulates how complementarity operates.

These four concepts answer two fundamental questions: \textbf{HOW} do utility components interact (complementarity $\gamma$), and \textbf{WHEN} does context modulate behavior ($\Psi$-dependence)? But to \emph{predict} behavior---not just explain it post hoc---we need to answer \textbf{six additional questions}:

\begin{itemize}[nosep]
    \item \textbf{WHO} has utility? (Thomas thinks as individual, not as family)
    \item \textbf{WHAT} is utility? (Thomas focuses on Physical, misses Social benefits)
    \item \textbf{WHERE} do parameters come from? (How do we estimate $\gamma$ for Thomas?)
    \item \textbf{AWARE}: How salient is the utility? (Thomas's long-term benefits are filtered)
    \item \textbf{READY}: How willing to act? (Thomas at $WAX = 0.6$, threshold $\theta = 0.7$)
    \item \textbf{STAGE}: Where in the change journey? (Thomas at Stage 1: Contemplation)
\end{itemize}

\noindent Thomas's Fitness Journey (above) illustrates what Anna's Dilemma could not: why \emph{knowing} the right answer and \emph{doing} the right thing are separated by an entire behavioral pipeline. This chapter develops that pipeline---the \textbf{10C CORE Framework}.

%------------------------------------------------------------------------------
% CHAPTER OVERVIEW
%------------------------------------------------------------------------------
\paragraph{Chapter Overview.}
\label{sec:10c-overview}
This chapter proceeds as follows:
\begin{itemize}
    \item Section \ref{sec:10c-questions}: The eight fundamental questions (overview)
    \item Section \ref{sec:10c-pipeline}: The four-stage pipeline from utility to action
    \item Section \ref{sec:10c-who}: CORE-WHO---The level hierarchy
    \item Section \ref{sec:10c-what}: CORE-WHAT---The FEPSDE dimensions
    \item Section \ref{sec:10c-how}: CORE-HOW---Complementarity
    \item Section \ref{sec:10c-when}: CORE-WHEN---The $\Psi$ context dimensions
    \item Section \ref{sec:10c-where}: CORE-WHERE---Calibration
    \item Section \ref{sec:10c-aware}: CORE-AWARE---The awareness filter
    \item Section \ref{sec:10c-ready}: CORE-READY---The willingness threshold
    \item Section \ref{sec:10c-stage}: CORE-STAGE---The behavioral change journey
    \item Section \ref{sec:10c-integration}: Integration---How the 10C connect
    \item Section \ref{sec:10c-summary}: Summary---The complete architecture
\end{itemize}

%==============================================================================
\subsection{The Eight Fundamental Questions}
\label{sec:10c-questions}
%==============================================================================

\begin{quote}
\textbf{Why do people often act differently than expected?}
\end{quote}

Standard economic theory predicts behavior from preferences and constraints.
Yet systematic deviations occur---not randomly, but in patterns that depend on
\emph{who} is deciding, \emph{what} they value, \emph{how} values interact,
\emph{when} context matters, \emph{where} parameters come from, how \emph{aware}
they are, and how \emph{ready} to act.

EBF answers this question through \textbf{eight fundamental questions}---the
\textbf{10C CORE Framework}.

%==============================================================================
\subsection{The Eight Fundamental Questions}
%==============================================================================

\begin{definition}[The 10C CORE Framework]
\label{def:10c-core}
Every behavioral prediction requires answering eight questions:

\begin{center}
\renewcommand{\arraystretch}{1.4}
\begin{tabular}{@{}clll@{}}
\toprule
\textbf{\#} & \textbf{CORE} & \textbf{Question} & \textbf{Output} \\
\midrule
1 & \textbf{WHO} & Who has utility? & Level $L \in \{0, 1, ..., \infty\}$ \\
2 & \textbf{WHAT} & What is utility? & Dimensions $d \in$ FEPSDE \\
3 & \textbf{HOW} & How do components interact? & Complementarity $\gamma$ \\
4 & \textbf{WHEN} & When does context matter? & Context $\Psi$ (8 dimensions) \\
5 & \textbf{WHERE} & Where do numbers come from? & Parameters $\Theta$ \\
6 & \textbf{AWARE} & How salient is the utility? & Awareness $A(\cdot)$ \\
7 & \textbf{READY} & How willing to act? & Willingness $WAX$, threshold $\theta$ \\
8 & \textbf{STAGE} & Where in the change journey? & Stage $S(t)$ \\
\bottomrule
\end{tabular}
\end{center}
\end{definition}

%==============================================================================
\subsection{The Four-Stage Pipeline}
\label{sec:10c-pipeline}
%==============================================================================

The 10C questions organize into a \textbf{four-stage pipeline} from potential
utility to observable action:

\begin{tcolorbox}[colback=blue!5!white, colframe=blue!75!black,
    title=\textbf{The EBF Pipeline: From Utility to Behavioral Change}]

\textbf{Stage 1: UTILITY DEFINITION} (5 questions)
\begin{equation}
U_{pot} = f(L, d, \gamma, \Psi, \Theta)
\end{equation}
\begin{itemize}[nosep]
\item \textbf{WHO} ($L$): At which level is utility evaluated?
\item \textbf{WHAT} ($d$): Which dimensions matter?
\item \textbf{HOW} ($\gamma$): How do dimensions interact?
\item \textbf{WHEN} ($\Psi$): How does context modulate?
\item \textbf{WHERE} ($\Theta$): What are the parameter values?
\end{itemize}

\vspace{0.3cm}
\textbf{Stage 2: AWARENESS FILTER} (6th question)
\begin{equation}
U_{eff} = A(\cdot) \times U_{pot}
\end{equation}
\begin{itemize}[nosep]
\item \textbf{AWARE} ($A$): How much of potential utility is psychologically salient?
\end{itemize}

\vspace{0.3cm}
\textbf{Stage 3: ACTION THRESHOLD} (7th question)
\begin{equation}
\text{Action} \Leftrightarrow WAX \geq \theta
\end{equation}
\begin{itemize}[nosep]
\item \textbf{READY} ($WAX, \theta$): Is willingness above the action threshold?
\end{itemize}

\vspace{0.3cm}
\textbf{Stage 4: BEHAVIORAL CHANGE JOURNEY} (8th question)
\begin{equation}
S(t+1) = f(S(t), A, WAX, \theta, \Psi)
\end{equation}
\begin{itemize}[nosep]
\item \textbf{STAGE} ($S(t)$): Where is the person in their behavioral change journey?
\end{itemize}

\end{tcolorbox}

\paragraph{Why Four Stages?}
This pipeline captures a fundamental truth about human behavior:

\begin{enumerate}
\item \textbf{Potential $\neq$ Effective:} People may have high potential utility
from an action but not be \emph{aware} of it. The Awareness filter ($A$)
determines what fraction of potential utility becomes psychologically effective.

\item \textbf{Effective $\neq$ Action:} People may be aware of utility but not
\emph{ready} to act. Barriers, friction, habits, and competing demands create
a threshold ($\theta$) that willingness ($WAX$) must exceed.

\item \textbf{Action $\neq$ Sustained Change:} A single action does not guarantee
behavioral change. The \textbf{STAGE} function captures where individuals are
in their change journey and how likely they are to maintain or regress.
\end{enumerate}

This explains why information campaigns often fail (Stage 2 bottleneck),
why intentions don't predict behavior (Stage 3 bottleneck), and why behavior
change programs often show relapse (Stage 4 dynamics).

%==============================================================================
\subsection{CORE-WHO: The Level Hierarchy}
\label{sec:10c-who}
%==============================================================================

\paragraph{The Innovation.}
Standard economics asks ``what maximizes utility?'' but assumes the answer is
obvious: the individual. EBF asks first: \textbf{WHO has utility?}

The answer is not singular. Utility is evaluated at multiple \textbf{aggregation
levels}, and which level is salient depends on context:

\begin{center}
\renewcommand{\arraystretch}{1.2}
\begin{tabular}{@{}cll@{}}
\toprule
\textbf{Level} & \textbf{Entity} & \textbf{Example Question} \\
\midrule
$L = 0$ & Individual (Self) & ``What's best for ME?'' \\
$L = 1$ & Dyad (Partner) & ``What's best for US (couple)?'' \\
$L = 2$ & Family/Household & ``What's best for OUR FAMILY?'' \\
$L = 3$ & Tribe/Community & ``What's best for OUR COMMUNITY?'' \\
$L = 4$ & Organization & ``What's best for THE COMPANY?'' \\
$L = 5$ & Nation & ``What's best for THE COUNTRY?'' \\
$L = 6$ & Civilization & ``What's best for HUMANITY?'' \\
$L = \infty$ & Meta & ``What's best for ALL BEINGS?'' \\
\bottomrule
\end{tabular}
\end{center}

\paragraph{Level Salience.}
Which level dominates is determined by the \textbf{level salience function}
$\alpha^L(\Psi)$, which depends on context:

\begin{equation}
\text{Total Utility} = \sum_{L=0}^{\infty} \alpha^L(\Psi) \cdot U^L
\end{equation}

where $\sum_L \alpha^L(\Psi) = 1$.

\textit{Cross-Reference:} The complete 9-level hierarchy and 41 axioms are
formalized in Appendix~AAA (CORE-WHO).

%==============================================================================
\subsection{CORE-WHAT: The FEPSDE Dimensions}
\label{sec:10c-what}
%==============================================================================

\paragraph{The Innovation.}
Standard economics reduces utility to consumption or income. EBF asks:
\textbf{WHAT is utility?}

The answer has six dimensions, each irreducible to the others:

\begin{center}
\renewcommand{\arraystretch}{1.2}
\begin{tabular}{@{}cll@{}}
\toprule
\textbf{Dim} & \textbf{Name} & \textbf{What It Captures} \\
\midrule
$F$ & Financial & Material wealth, income, economic security \\
$E$ & Emotional & Psychological wellbeing, satisfaction, meaning \\
$P$ & Physical & Health, energy, bodily integrity, longevity \\
$S$ & Social & Relationships, belonging, trust, status \\
$D$ & Digital & Connectivity, access, digital participation \\
$Eco$ & Ecological & Environmental quality, sustainability \\
\bottomrule
\end{tabular}
\end{center}

Combined with the three utility categories (INU, KNU, IDN), four time horizons,
and gain/pain valence, this yields the \textbf{144-component structure}:

\begin{equation}
3 \text{ (INU/KNU/IDN)} \times 6 \text{ (FEPSDE)} \times 4 \text{ (time)} \times 2 \text{ (valence)} = 144
\end{equation}

\textit{Cross-Reference:} The complete matrix and axioms are formalized in
Appendix~C (CORE-WHAT).

%==============================================================================
\subsection{CORE-HOW: Complementarity}
\label{sec:10c-how}
%==============================================================================

\paragraph{The Innovation.}
Standard economics assumes separability: your choices don't affect my optimal
choice (except through prices). EBF asks: \textbf{HOW do components interact?}

The answer is \textbf{complementarity}---the cross-partial derivative:

\begin{equation}
\gamma_{ij} = \frac{\partial^2 U}{\partial x_i \partial x_j}
\end{equation}

\begin{itemize}
\item $\gamma_{ij} > 0$: Complements---more of $i$ increases marginal value of $j$
\item $\gamma_{ij} < 0$: Substitutes---more of $i$ decreases marginal value of $j$
\item $\gamma_{ij} = 0$: Independent---the Arrow-Debreu case
\end{itemize}

Complementarity creates \textbf{non-separability}, which generates:
\begin{itemize}[nosep]
\item Multiple equilibria
\item Path dependence
\item Regime shifts
\item The possibility of coordination failure
\end{itemize}

\textit{Cross-Reference:} The complete complementarity framework is developed
in Chapter~5 and formalized in Appendix~B (CORE-HOW).

%==============================================================================
\subsection{CORE-WHEN: The $\Psi$ Context Dimensions}
\label{sec:10c-when}
%==============================================================================

\paragraph{The Innovation.}
Standard economics treats context as exogenous background. EBF asks:
\textbf{WHEN does context matter?}

The answer: \textbf{always}, through eight systematic dimensions:

\begin{center}
\small
\renewcommand{\arraystretch}{1.2}
\begin{tabular}{@{}clll@{}}
\toprule
\textbf{Symbol} & \textbf{Dimension} & \textbf{What It Captures} \\
\midrule
$\Psi_I$ & Institutional & Rules, enforcement, property rights \\
$\Psi_S$ & Social & Trust, norms, social capital \\
$\Psi_C$ & Cognitive & Information processing, numeracy \\
$\Psi_K$ & Informational & Transparency, signal quality \\
$\Psi_E$ & Economic & Development, market depth \\
$\Psi_T$ & Temporal & Time horizons, stability \\
$\Psi_M$ & Market Scope & Geographic reach \\
$\Psi_F$ & Factor Flexibility & Adjustment capacity \\
\bottomrule
\end{tabular}
\end{center}

Context is not merely background---it \textbf{modulates} every parameter:
loss aversion, discount rates, reference points, complementarities.

\textit{Cross-Reference:} The complete context framework is developed in
Chapter~9 and formalized in Appendix~V (CORE-WHEN).

%==============================================================================
\subsection{CORE-WHERE: Calibration}
\label{sec:10c-where}
%==============================================================================

\paragraph{The Innovation.}
Most theories leave parameters unspecified. EBF asks: \textbf{WHERE do
the numbers come from?}

The answer involves systematic calibration using:
\begin{enumerate}
\item Experimental data (lab and field)
\item Survey measures (stated preferences)
\item Revealed preferences (observed choices)
\item \textbf{LLM Monte Carlo} calibration (new method, see Chapter~4x)
\end{enumerate}

All parameters carry \textbf{epistemic status tags}:
\begin{itemize}[nosep]
\item \texttt{[EMP]}: Empirically validated
\item \texttt{[THR]}: Theoretically derived
\item \texttt{[LLM]}: LLM Monte Carlo estimated
\item \texttt{[ILL]}: Illustrative
\item \texttt{[HYP]}: Hypothetical
\end{itemize}

\textit{Cross-Reference:} The complete estimation methodology is formalized
in Appendix~BBB (CORE-WHERE).

%==============================================================================
\subsection{CORE-AWARE: The Awareness Filter}
\label{sec:10c-aware}
%==============================================================================

\paragraph{The Innovation.}
Standard economics assumes people know their utility function. EBF asks:
\textbf{How AWARE are people of their utility?}

The Awareness function $A(\cdot)$ filters potential utility to effective utility:

\begin{equation}
U_{eff} = A(U_{pot}, \Psi, \text{attention}, \text{salience})
\end{equation}

Key mechanisms:
\begin{itemize}[nosep]
\item \textbf{Attention:} Limited cognitive bandwidth
\item \textbf{Salience:} Some dimensions more visible than others
\item \textbf{Moment of Truth:} When awareness peaks
\item \textbf{Motivated reasoning:} Beliefs shaped by desires
\end{itemize}

The \textbf{Transformation Equation}:
\begin{equation}
U_{eff}^d = s^d(\Psi) \cdot \phi^d(U_{pot}^d)
\end{equation}

where $s^d$ is dimension-specific salience and $\phi^d$ is the attention function.

\textit{Cross-Reference:} The complete awareness framework with 104 axioms is
developed in Chapter~11 and formalized in Appendix~AU (CORE-AWARE).

%==============================================================================
\subsection{CORE-READY: The Willingness Threshold}
\label{sec:10c-ready}
%==============================================================================

\paragraph{The Innovation.}
Standard economics assumes that if utility is positive, action follows.
EBF asks: \textbf{How READY are people to act?}

The Willingness-to-Act index ($WAX$) must exceed a threshold ($\theta$):

\begin{equation}
\boxed{\text{Action} \Leftrightarrow WAX \geq \theta}
\end{equation}

$WAX$ integrates:
\begin{itemize}[nosep]
\item Effective utility $U_{eff}$
\item Thinking style (System 1 vs. System 2)
\item Behavioral change stage (BCJ)
\item Barriers and friction
\end{itemize}

The threshold $\theta$ depends on:
\begin{itemize}[nosep]
\item Default options
\item Social proof
\item Commitment devices
\item Implementation intentions
\end{itemize}

\textit{Cross-Reference:} The complete willingness framework with 99+ axioms is
developed in Chapter~12 and formalized in Appendix~AV (CORE-READY).

%==============================================================================
\subsection{CORE-STAGE: The Behavioral Change Journey}
\label{sec:10c-stage}
%==============================================================================

\paragraph{The Innovation.}
Standard economics assumes behavior change is instantaneous: once utility is
positive, action follows immediately and permanently. EBF asks:
\textbf{Where in the CHANGE journey is the person?}

The Stage function $S(t)$ tracks position through behavioral change:

\begin{equation}
S(t) \in \{0, 1, 2, 3, 4, 5\}
\end{equation}

\begin{center}
\renewcommand{\arraystretch}{1.2}
\begin{tabular}{@{}cll@{}}
\toprule
\textbf{Stage} & \textbf{Name} & \textbf{Characteristic} \\
\midrule
$S = 0$ & Pre-contemplation & Not aware of need for change \\
$S = 1$ & Contemplation & Considering change \\
$S = 2$ & Preparation & Planning to change \\
$S = 3$ & Action & Actively changing \\
$S = 4$ & Maintenance & Sustaining change \\
$S = 5$ & Termination & Change fully integrated \\
\bottomrule
\end{tabular}
\end{center}

The \textbf{Stage Dynamics} equation:
\begin{equation}
\frac{dS}{dt} = f(A, WAX, \theta, \tau, \rho, \beta, \Psi)
\end{equation}

where $\tau$ is transition probability, $\rho$ is relapse rate, and $\beta$
is social contagion strength.

\textit{Cross-Reference:} The complete behavioral change journey with 100+ axioms is
developed in Chapter~13 and formalized in Appendix~AW (CORE-STAGE).

%==============================================================================
\subsection{Integration: How the 10C Connect}
\label{sec:10c-integration}
%==============================================================================

The eight COREs are not independent---they form an integrated architecture:

\begin{tcolorbox}[colback=green!5!white, colframe=green!75!black,
    title=\textbf{The 10C Integration Structure}]

\begin{center}
\begin{tabular}{@{}l@{}}
\textbf{WHO} determines at which level utility is evaluated \\
$\downarrow$ \\
\textbf{WHAT} specifies the dimensions at that level \\
$\downarrow$ \\
\textbf{HOW} captures interactions between dimensions \\
$\downarrow$ \\
\textbf{WHEN} modulates all parameters via context $\Psi$ \\
$\downarrow$ \\
\textbf{WHERE} calibrates the numerical values \\
$\downarrow$ \\
\textbf{AWARE} filters potential to effective utility \\
$\downarrow$ \\
\textbf{READY} determines if effective utility triggers action \\
$\downarrow$ \\
\textbf{STAGE} tracks position in behavioral change journey \\
\end{tabular}
\end{center}

\textbf{Key Insight:} Omitting any CORE leads to prediction failure:
\begin{itemize}[nosep]
\item Miss WHO $\rightarrow$ wrong aggregation level
\item Miss WHAT $\rightarrow$ wrong dimensions
\item Miss HOW $\rightarrow$ miss interaction effects
\item Miss WHEN $\rightarrow$ context-blind prediction
\item Miss WHERE $\rightarrow$ ungrounded parameters
\item Miss AWARE $\rightarrow$ assume perfect knowledge
\item Miss READY $\rightarrow$ assume intention = action
\item Miss STAGE $\rightarrow$ assume one-shot behavior
\end{itemize}

\end{tcolorbox}

%==============================================================================
\subsection{Connecting to the Four Core Concepts}
\label{sec:10c-connect}
%==============================================================================

The 10C CORE Framework \emph{extends} the four concepts introduced in
Chapter~1 (Section~\ref{sec:introduction}):

\begin{center}
\renewcommand{\arraystretch}{1.3}
\begin{tabular}{@{}lll@{}}
\toprule
\textbf{Original Concept} & \textbf{10C Connection} & \textbf{What 10C Adds} \\
\midrule
Complementarity $C_{ij}$ & CORE-HOW ($\gamma$) & Level-specific $\gamma^L$ \\
Context $\Psi$ & CORE-WHEN & Systematic 8-dimension structure \\
Reference $C^*(\Psi)$ & All 10C & $C^*$ now includes WHO, WHAT \\
Coherence $K$ & All 10C & $K$ measures fit across all 8 \\
Change dynamics & CORE-STAGE & Temporal $S(t)$ trajectory \\
\bottomrule
\end{tabular}
\end{center}

The 10C Framework answers the question: \textbf{What determines $C^*$?}

\begin{equation}
C^*(\Psi, t) = f\underbrace{(L}_{\text{WHO}}, \underbrace{d}_{\text{WHAT}}, \underbrace{\gamma}_{\text{HOW}}, \underbrace{\Psi}_{\text{WHEN}}, \underbrace{\Theta}_{\text{WHERE}}, \underbrace{A}_{\text{AWARE}}, \underbrace{WAX, \theta}_{\text{READY}}, \underbrace{S(t)}_{\text{STAGE}})
\end{equation}

%==============================================================================
\subsection{Summary: The Complete Architecture}
\label{sec:10c-summary}
%==============================================================================

\begin{tcolorbox}[colback=gray!5!white, colframe=gray!75!black,
    title=\textbf{EBF: The 10C CORE Framework}]

\textbf{The Central Equation:}
\begin{equation}
\text{Sustained Action} \Leftrightarrow WAX\Big(A\big(U_{pot}(L, d, \gamma, \Psi, \Theta)\big)\Big) \geq \theta(\Psi) \text{ AND } S(t) \text{ progresses}
\end{equation}

\textbf{In Words:}
\begin{quote}
Sustained behavioral change occurs when willingness-to-act (based on effective utility, which is
awareness-filtered potential utility defined over levels, dimensions, and
complementarities in context) exceeds the action threshold AND the individual
progresses through the behavioral change stages.
\end{quote}

\textbf{The Eight Questions:}
\begin{enumerate}
\item \textbf{WHO} has utility? $\rightarrow$ Level $L$
\item \textbf{WHAT} is utility? $\rightarrow$ Dimensions $d$
\item \textbf{HOW} do they interact? $\rightarrow$ Complementarity $\gamma$
\item \textbf{WHEN} does context matter? $\rightarrow$ Context $\Psi$
\item \textbf{WHERE} do numbers come from? $\rightarrow$ Parameters $\Theta$
\item How \textbf{AWARE}? $\rightarrow$ Awareness filter $A(\cdot)$
\item How \textbf{READY}? $\rightarrow$ Willingness $WAX$, threshold $\theta$
\item What \textbf{STAGE}? $\rightarrow$ Stage position $S(t)$
\end{enumerate}

\end{tcolorbox}

The remainder of this document develops each CORE in detail, demonstrates
empirical applications, and derives policy implications.

%------------------------------------------------------------------------------
% READING PATH BOX
%------------------------------------------------------------------------------
\begin{tcolorbox}[colback=green!5!white, colframe=green!75!black,
    title=\textbf{Reading Paths from Chapter 1b}]
\small

\textbf{Continue reading:}
\begin{itemize}[nosep]
    \item \textbf{Chapter 2:} Rationality concepts in classical economics
    \item \textbf{Chapter 5:} Complementarity (CORE-HOW in depth)
\end{itemize}

\medskip
\textbf{Deep dive into individual COREs:}
\begin{itemize}[nosep]
    \item \textbf{WHO:} Appendix WHO (CORE-WHO) + Chapter 10.8
    \item \textbf{WHAT:} Appendix WAT (CORE-WHAT) + Chapter 10.2
    \item \textbf{HOW:} Appendix HOW (CORE-HOW) + Chapter 5
    \item \textbf{WHEN:} Appendix WEN (CORE-WHEN) + Chapter 9
    \item \textbf{WHERE:} Appendix EST (CORE-WHERE) + Chapter 4x
    \item \textbf{AWARE:} Appendix AWA (CORE-AWARE) + Chapter 11
    \item \textbf{READY:} Appendix REA (CORE-READY) + Chapter 12
    \item \textbf{STAGE:} Appendix STA (CORE-STAGE) + Chapter 13
\end{itemize}

\medskip
\textbf{Application-focused readers:}
\begin{itemize}[nosep]
    \item Skip to Chapter 14 (Behavioral Change Segments) for practical applications
    \item See Chapter 15 (WEC Synthesis) for the integrated prediction framework
\end{itemize}
\end{tcolorbox}

%%%%%%%%%%%%%%%%%%%%%%%%%%%%%%%%%%%%%%%%%%%%%%%%%%%%%%%%%%%%%%%%%%%%%%%%%%%%%%%
% END OF CHAPTER 1b: THE 10C CORE ARCHITECTURE
%%%%%%%%%%%%%%%%%%%%%%%%%%%%%%%%%%%%%%%%%%%%%%%%%%%%%%%%%%%%%%%%%%%%%%%%%%%%%%%
